% !TeX root = ../main.tex
\section{Introduction}
\label{sec:intro}

Shear walls are the principal lateral-force-resisting components in many high-rise residential buildings, particularly in seismic regions \cite{GB500112010Code2016,qianDesignTallBuilding2018}. The arrangement of shear walls influences stiffness distribution, load paths, and construction cost, and it must simultaneously respect architectural constraints such as functional zoning, circulation, and openings. In preliminary design, these competing requirements are commonly reconciled through iterative manual layout adjustment and repeated coordination between architectural and structural teams \cite{Steiner_2016}. Such a workflow is labor-intensive and tends to restrict systematic exploration of alternative layouts.

Despite extensive design experience and code provisions, routine shear wall layout development remains inefficient. In practice, engineers still rely on implicit, experience-based rules to translate architectural drawings into structurally rational wall candidates, and these rules are difficult to formalize or reuse across projects. Once initial candidates are proposed, cross-discipline conflict resolution typically requires repeated coordination between architectural and structural teams, which slows early-stage decision making \cite{heGenerativeAIBIMAutomatic2025}. This workflow also narrows the range of schemes that can be explored within realistic time constraints, so final layouts may be acceptable locally but suboptimal at the global plan level. These issues motivate automated methods that can encode architectural constraints and generate feasible wall layouts with explicit control over design intent \cite{liaoAutomatedStructuralDesign2021a}.

Recent data-driven methods have treated shear wall generation as image-to-image translation, using convolutional networks, GAN-based models, and diffusion models to map floor-plan images to structural layouts \cite{pizarroUseConvolutionalNetworks2021,zhouStructDiffusionEndendIntelligent2024}. While pixel-based representations can capture complex spatial patterns, they are inefficient for structural elements because thin wall lines occupy only a small fraction of the image grid, topological relationships remain implicit, and raster outputs still require vectorization for downstream engineering analysis \cite{wangHighResolutionImageSynthesis2018,feiIntegratedSchematicDesign2022}.

To address the issue of implicit topology, graph neural networks (GNNs) have been adopted to model structural systems more explicitly \cite{Bronstein_2017,Hamilton_2020}. A common formulation builds graphs from geometric components, with nodes as wall intersections and edges as wall segments \cite{zhaoIntelligentDesignShear2023b}. Based on it, structural layout prediction is formulated as either a classification or regression task. This representation preserves connectivity but remains low-level, as individual segments carry limited information about the architectural spaces they bound. As a result, room-level constraints (e.g., functional requirements, openings, and circulation) are not directly expressed. A room-level representation is more aligned with the physical and architectural constraints governing shear wall placement. In typical residential plans, feasible shear wall candidates are largely restricted to room boundaries, and many practical constraints are naturally described at the room level rather than at the level of wall segments. Modeling rooms as graph nodes and room adjacencies as edges therefore provides an explicit interface between architectural semantics and structural layout generation \cite{nauataHouseGANRelationalGenerative2020}. In addition, the number of rooms per plan is usually far smaller than the number of component wall segments, which reduces graph scale and can improve computational efficiency. Related room-graph and room-semantics studies in architectural layout generation also support the effectiveness of this abstraction \cite{huGraph2PlanLearningFloorplan2020}.

Beyond representation, structural layouts are not determined by architectural plans alone; they also depend on design conditions such as seismic intensity and building height \cite{zhaoDesignconditioninformedShearWall2023d}. A model that ignores these conditions may converge to statistically averaged layouts that are not useful for engineering decisions. At the same time, robust conditional learning is difficult because each floor plan is usually paired with only one finalized layout in real projects. Therefore, this study targets two coupled challenges: introducing room-level semantics into graph-based generation and maintaining condition responsiveness under sparse paired supervision.

To address these challenges, we make three contributions with distinct roles in the framework. We first introduce a room-level graph formulation that treats rooms and room adjacencies as the core prediction units, so the model can reason directly over boundary-feasible locations while reducing graph scale relative to component-level graphs. We then adopt FiLM-based conditional modulation inside the GNN backbone to keep condition signals active through deep message passing and improve adherence to specified density regimes. Finally, we design a dual-stream training strategy that combines paired supervision with cross-condition density regularization, which provides condition-related gradients even when multi-condition labels are unavailable for the same plan. Together, these designs connect architectural semantics, controllable generation, and practical data constraints in a unified pipeline for preliminary shear wall layout generation.
